\section{No Interior Mixed Term}

Completion is often where forbidden mixtures try to return. The finite
chapter had already split the geometric and gauge carriers. A finite
certificate could read the geometric face or the gauge face, but it could
not form an interior direct sum that smuggled in a scalar mixed term. The
Sobolev--Hilbert lift preserves that split. Therefore the completed grammar
must also forbid the interior product that would undo it.

The issue is not whether geometry and gauge are connected. The previous
chapter established the connection sharply: boundary geometry generated the
gauge. The issue is where the connection may be read. Since the gauge arose
because the interior could not settle the charged residue on its own, the
completed interior may not later claim a scalar certificate that multiplies
geometric and gauge residues as if both were freely present there.

Compatibility supplies the rule. A continuum reading is compatible only if
it preserves the obligations that finite refinement forced. It must preserve
extensional countability, successor residue, the normed completion, and the
geometry/gauge split. A smooth-looking expression that violates one of those
obligations is not a deeper reading. It is an inadmissible presentation.

Thus the interior mixed term cancels:
\[
  \langle g,q\rangle_{\mathrm{int}} = 0 .
\]
This is not a claim that the symbols \(g\) and \(q\) have no relation. It is
the statement that their interior scalar cross-reading is forbidden by
compatibility. The geometric channel \(g\) and the gauge channel \(q\) may
meet at the boundary process that generated the gauge, but the completed
interior cannot multiply them into a hidden scalar debt.

The cancellation has a grammatical source. The split lifted as
orthogonality. An interior mixed scalar would identify what orthogonality was
forced to keep apart. It would let the geometric channel pay gauge debt or
the gauge channel pay geometric debt inside the interior. That would erase
the boundary origin of the gauge and turn the forced split into a notation.
Compatibility forbids that erasure.

The interior is therefore poorer than the completed notation might suggest.
It can carry completed geometric readings. It can carry completed gauge
readings. It can carry the normed residue histories that led to both. But it
cannot carry an independent scalar product between the channels. The missing
term is not a gap. It is the discipline that keeps the completed reading
faithful to the finite one.

A Casimir-style boundary imbalance gives a useful illustration. When a
boundary restricts admissible interior refinements, the repair cannot be
made by inventing modes inside the restricted region. The repair is read at
the boundary where the imbalance is exposed. The example does not prove the
chapter's claim, but it shows the same grammatical shape: an interior
deficit is not healed by adding unearned interior structure.

The holographic illustration says the same thing at a larger scale. If every
admissible interior refinement must be reconciled through the boundary, then
volume is not a storehouse of independent record. Interior presentations are
smooth shadows of boundary bookkeeping. The present chapter is narrower: it
does not need a general scaling law. It needs only the rule that the
geometry/gauge mixed scalar has no independent interior certificate.

This also clarifies the role of compatibility in the continuum. Compatibility
is not passive agreement between completed objects. It is the active refusal
to let completion forget how the objects were earned. If a completed reading
would require an interior mixed term, then it fails compatibility with the
finite split. The grammar cancels the term before it can become a false
reconciliation.

The residue, however, has not vanished from the story. Removing the interior
mixed scalar does not remove the terminal debt that the geometric and gauge
readings jointly expose. It only says where that debt cannot be paid. The
residue must still be read, and the only remaining licensed place is the
boundary trace. The next section therefore forces radiation to the boundary.
